\subsection{Private-market structures, performance and the investment process}

\begin{qfmethod}[Diversification benefits of private capital: low reported correlations, smoothing and unsmoothing, cycle-dependent strategies, and the opposing views]\label{meth:private-capital-diversification}
\qffield{Statement} Private equity and debt add exposures that public markets lack: companies and loans that are not listed, control-oriented strategies, and strategies that do well at different points of the cycle. Demaria and co-authors argue that venture capital does best when invested at the end of a recession, buyouts when assets are bought just after a peak, and distressed debt when equity prices and growth fall; distressed debt is ``anti-correlated'' with venture capital \citep[pp.~162--163]{demaria-private-markets-2021}.

\emph{Measured correlation.} Reported returns are appraisals, which lag and smooth true values. With an AR(1) smoother $r^*_t=(1-\phi)r_t+\phi r^*_{t-1}$, the reported variance is $(1-\phi)/(1+\phi)$ of the true variance, and contemporaneous betas and correlations shrink in proportion (Derivation). Ang concludes that reported illiquid-asset returns ``are not returns''; volatilities, correlations and betas from them are too low. The Geltner and Ross--Zisler correction inverts the smoother \citep[pp.~416--421]{s7-ang}:
\[ r_t=\frac{r^*_t-\phi\,r^*_{t-1}}{1-\phi}. \]
Getmansky, Lo and Makarov model hedge-fund serial correlation the same way, as a sign of illiquidity exposure \citep[NBER WP 9571, abstract]{wide-getmansky-lo-makarov}.

\emph{The opposing view.} Demaria and co-authors call de-smoothing buyout returns a ``myth''. Appraisals of control stakes reflect rights, warranties and premiums that no listed price captures, so more frequent proxies add noise, and fund-of-funds return distributions show low loss rates \citep[pp.~160--166]{demaria-private-markets-2021}. The evidence on cash-flow performance shows buyouts beating public markets on average net of fees \citep[abstract]{pe-harris-jenkinson-kaplan-2014}, but not that their risk is low.

\qffield{Practice} \emph{Smoothing simulated} (original, 40,000 quarters, seed fixed). True quarterly returns with volatility 10 percent and market beta 1.0 are reported through a smoother with $\phi=0.6$:
\begin{itemize}
\item reported volatility is 5 percent and contemporaneous beta 0.40, so the reported Sharpe ratio doubles;
\item the estimated autocorrelation is 0.60, and unsmoothing restores volatility and beta to within 0.1 percent.
\end{itemize}
An allocator who uses the reported figures would see a diversifier with half its true risk. The cash-flow PME of a fund is unaffected by smoothing, because it uses only calls and distributions (\hyperref[meth:private-fund-multiples-pme]{multiples and PME}).

\qffield{Conditions} The AR(1) smoother is a model; real appraisal lags vary by asset and manager. The cycle claims are the authors' interpretation of patterns, not tests.

\qffield{Pitfalls} Optimising an allocation on appraisal-based volatilities and correlations. Concluding that unsmoothed returns are the truth rather than another estimate. Treating cycle diversification as available to an investor who commits pro-cyclically.

\qffield{See also} \hyperref[meth:private-public-market-comparison]{private against public markets}; \hyperref[meth:alternatives-allocation-approaches]{allocation}.

\end{qfmethod}

\begin{qfmethod}[Private against public investments and markets: contrasting features, and comparing their risk and return within a strategic asset allocation]\label{meth:private-public-market-comparison}
\qffield{Statement} \emph{Contrasts} \citep[p.~594, Table~18.1]{s7-ang}:
\begin{center}\small
\begin{tabular}{@{}p{0.2\linewidth}p{0.34\linewidth}p{0.38\linewidth}@{}}
\toprule
 & Public markets & Private markets \\
\midrule
Trading & centralised, liquid & over the counter, illiquid, search frictions \\
Transaction costs & small & large; secondary discounts of 15--30 percent, up to 50--70 percent in crises \\
Valuation & objective, continuous & appraised quarterly, subjective \\
Horizon & immediate & about ten years \\
Contracts & standard & complex, negotiated \\
Information & public disclosure & private, asymmetric \\
Owner's role & mostly passive & control and monitoring \\
\bottomrule
\end{tabular}
\end{center}
Secondary discounts: \citep[p.~63]{demaria-private-markets-2021}.

\emph{Risk and return.} Reported private returns overstate the reward and understate risk because of survivorship, smoothing and selection \citep[pp.~416--426]{s7-ang}. Ang argues that, adjusted for these and for leverage, the average private equity fund provides no excess return and that ``there is no market index'' an investor can buy, so idiosyncratic manager risk dominates (pp.~427--428). Cash-flow evidence from Burgiss shows average US buyout funds beating the S\&P~500 by more than 3 percent a year net of fees \citep[abstract]{pe-harris-jenkinson-kaplan-2014}. The two views differ mainly on the risk adjustment and on what an investor without top-quartile access can expect.

\emph{For a strategic allocation} (Interpretation): use unsmoothed volatilities and betas, treat the private premium as compensation for illiquidity, leverage and manager risk rather than as free return, and set the weight subject to liquidity needs and the ability to keep commitments (\hyperref[meth:alternatives-allocation-approaches]{allocation approaches}).

\qffield{Practice} \emph{Reported against unsmoothed inputs} (original, mean--variance with risk aversion 3, no short positions). Expected returns are 4.5 percent for bonds, 7.5 percent for public equity and an assumed 10.5 percent for private equity.
\begin{itemize}
\item On reported figures (volatility 9 percent, correlation with public equity 0.45, beta 0.24) private equity has a Sharpe ratio of 0.83 and the optimiser puts everything in it; that portfolio's true volatility is 26 percent.
\item On unsmoothed figures (volatility 26 percent, correlation 0.85, beta 1.30) its Sharpe ratio is 0.29 against public equity's 0.26. The weights become 67 percent bonds, 1 percent public equity and 32 percent private equity, with volatility 9.7 percent.
\end{itemize}
\emph{Decision.} Even on unsmoothed inputs the optimiser treats private equity as a leveraged substitute for public equity, so a liquidity limit, not the optimiser, should bind. A pension with predictable outflows might hold 10--15 percent, implemented over several vintages.

\qffield{Conditions} The 10.5 percent expectation is an assumption for illustration, not a forecast; the historical outperformance depends on vintage, data source and benchmark.

\qffield{Pitfalls} Optimising with appraisal-based statistics. Comparing a fund's IRR with an index's time-weighted return. Assuming the median investor earns the average reported return.

\qffield{See also} \hyperref[meth:private-capital-diversification]{diversification}; \hyperref[meth:private-fund-multiples-pme]{performance measures}; \hyperref[def:private-equity-features]{private equity}.

\end{qfmethod}

\begin{qfmethod}[Private-market performance measures: DPI, RVPI and TVPI, the IRR and the J-curve, the public market equivalent, and why private and public performance differ]\label{meth:private-fund-multiples-pme}
\qffield{Statement} With paid-in capital $\mathrm{PIC}$ (calls, including fees), cumulative distributions $D$ and residual net asset value $\mathrm{NAV}$ \citep[p.~19]{ippolito-private-capital-2020}:
\[ \mathrm{DPI}=\frac{D}{\mathrm{PIC}},\qquad \mathrm{RVPI}=\frac{\mathrm{NAV}}{\mathrm{PIC}},\qquad \mathrm{TVPI}=\mathrm{DPI}+\mathrm{RVPI}. \]
DPI is realised cash; DPI above one means the fund has returned what was paid in. TVPI adds unrealised value, so LPs rely on it early in a fund's life and on DPI later.

\emph{IRR and the J-curve.} The since-inception IRR of calls, distributions and the current NAV is negative early, because fees and investments at cost precede gains, and rises as exits occur: the J-curve (pp.~16--18). The IRR measures return on capital invested, not on capital committed; it rewards early distributions and can be raised by pooling funds \citep[pp.~598--600]{s7-ang}; \citep[pp.~68--69]{demaria-private-markets-2021}.

\emph{Public market equivalent.} Kaplan and Schoar discount the fund's distributions and calls at the realised total return of a public index, and divide the first sum by the second; a PME above one means the fund beat the index net of fees \citep[NBER WP 9807, pp.~6, 8]{pe-kaplan-schoar-2005}. Their main sample is of largely liquidated funds, whose unrealised value is under 10 percent of commitments (p.~6). For a fund still holding investments, the residual NAV at the valuation date $T$ is treated as a final distribution \citep[pp.~10--12]{pe-harris-jenkinson-kaplan-2014}:
\[ \mathrm{PME}_T=\frac{\sum_{t\le T} D_t/I_t+\mathrm{NAV}_T/I_T}{\sum_{t\le T} C_t/I_t}. \]
Here $C_t$ and $D_t$ are the LP's calls and distributions net of fees and carry, $\mathrm{NAV}_T$ is the LP's share of residual value after accrued carry, and $I_t$ is the total-return index level on the same dates. For a liquidated fund $\mathrm{NAV}_T=0$ and the formula is Kaplan and Schoar's. The measure assumes the fund's beta is one; if it is higher, PME overstates risk-adjusted performance \citep[pp.~8--9]{pe-kaplan-schoar-2005}.

\emph{Excess-return versions.} The IRR of the index-deflated flows $C_t/I_t$ and $D_t/I_t$ (with $\mathrm{NAV}_T/I_T$) is a money-weighted excess return over the index. It is not the fund's IRR less the index's annualised return: with a constant index return $g$ it equals $(1+\mathrm{IRR})/(1+g)-1$, and with a varying index it depends on the timing of the flows (Derivation). Harris, Jenkinson and Kaplan describe a different measure, the Long--Nickels excess return, which subtracts the IRR of an index investment with the same calls (p.~11).

\emph{Why private and public performance differ.} Private returns are money-weighted and depend on the GP's timing of calls and exits; public returns are time-weighted. NAVs are appraisals that lag; multiples ignore time; and fees are paid on committed capital (Interpretation, from the sources above).

\qffield{Practice} \emph{Ashby Fund I} (original, commitment 100, four buyouts, European waterfall):
\begin{center}\small
\begin{tabular}{@{}lrrrrrr@{}}
\toprule
Year & 1 & 2 & 3 & 4 & 5 & 7 \\
\midrule
DPI & 0.00 & 0.00 & 0.00 & 0.67 & 1.11 & 1.52 \\
RVPI & 0.93 & 1.03 & 1.15 & 0.64 & 0.30 & 0.00 \\
TVPI & 0.93 & 1.03 & 1.15 & 1.31 & 1.41 & 1.52 \\
IRR since inception & $-14.8\%$ & 3.0\% & 8.0\% & 10.1\% & 11.0\% & 11.7\% \\
\bottomrule
\end{tabular}
\end{center}
\begin{itemize}
\item DPI passes one in year 5 and dips to 1.10 in year 6 as fees are still called.
\item With an index that returns 10, $-5$, 12, 8, 6, 15, $-10$ and 9 percent in years 1 to 8, the PME at liquidation is 1.23. The IRR of the index-deflated cash flows is 5.67 percent a year, a money-weighted margin over the index.
\item At the end of year 4 the LP holds NAV worth 0.64 of paid-in capital, and the interim PME including it is 1.10. Leaving the NAV out gives 0.56, as if the unrealised value were lost.
\item \emph{Benchmark case}: an LP contributes 100, receives nothing and holds NAV of 110 a year later while the index rises 10 percent. $\mathrm{PME}=(110/1.10)/(100/1.00)=1$: the fund matched the index. The distributions-only formula would give 0.
\end{itemize}

\qffield{Conditions} RVPI, interim IRRs and interim PMEs depend on the GP's valuations (IPEV guidelines in Europe). The PME uses one index, measured with dividends reinvested on the cash-flow dates; a small-cap or sector index may be the better comparator.

\qffield{Pitfalls} Ranking funds of different ages by TVPI. Reading a high interim IRR from a quick early exit as skill. Treating PME as a risk adjustment for leverage. Computing an interim PME without the residual NAV.

\qffield{See also} \hyperref[meth:gp-alignment-net-returns]{net returns and fees}; \hyperref[meth:alternative-returns-after-fees]{returns after fees}; \hyperref[meth:gips-return-calculations]{GIPS returns}.

\end{qfmethod}

\begin{qfmethod}[The general partner's role, alignment of interests with investors, and private-fund returns from the investor's side after management fees and carried interest]\label{meth:gp-alignment-net-returns}
\qffield{Statement} \emph{Roles.} The GP raises the fund, sources, analyses, invests in and monitors portfolio companies, implements value-creation actions, and exits; it decides alone, so LPs are passive \citep[pp.~13--15]{ippolito-private-capital-2020}. It also values the portfolio and reports to LPs.

\emph{Alignment and its limits.}
\begin{itemize}
\item Carried interest, a hurdle and a clawback tie the GP's pay to realised profits; the GP's own commitment gives it a stake in losses (pp.~14--16).
\item But management fees on committed capital, two-thirds of expected GP revenue \citep[p.~35]{metrick-yasuda-survey-2011}, reward size and delay. Transaction and monitoring fees reward activity, the IRR rewards early exits, and the carry option rewards risk \citep[pp.~610--614]{s7-ang}.
\item Remedies include fee offsets, a whole-of-fund waterfall, clawbacks, key-person and removal clauses, an LP advisory committee, and side letters and most-favoured-nation terms for some investors (pp.~611--612); limits on GP activity outside the fund are negative covenants (Metrick and Yasuda, p.~33).
\end{itemize}

\emph{Investor returns.} LP net returns follow from dated calls (investments plus fees) and distributions after carry. With a whole-of-fund waterfall and full catch-up, carry equals the carry rate times total fund profit once the hurdle is met (Derivation).

\qffield{Practice} \emph{Ashby Fund I} (original). Commitment 100; 2 percent fee on commitment for five years, then 1.5 percent of cost still held; 8 percent compounding hurdle, full catch-up, 20 percent carry. Deals: A (25 in year 0, sold for 60 in year 4), B (25 in year 1, 40 in year 5), C (20 in year 2, 50 in year 7), D (10 in year 3, written off in year 6).
\begin{itemize}
\item Paid-in capital 90.75, of which 10.75 fees. Proceeds 150, fund profit 59.25.
\item LPs receive 60, 40 and 38.15; the GP's carry of 11.85 is all paid in year 7, after the catch-up.
\item Gross deal IRR 16.87 percent and multiple 1.88; LP net IRR 11.71 percent and TVPI 1.52: fees and carry cost 5.2 points a year.
\item Deal by deal, the GP would take 6.33 in year 4, 1.42 in year 5 and 5.46 in year 7: 13.21, of which 1.36 must be clawed back because deal D's loss was never netted.
\end{itemize}

\qffield{Conditions} Allocating fees to deals by cost, as in the deal-by-deal case, is one convention among several. The clawback is only as good as the GP's ability to repay, often net of tax.

\qffield{Pitfalls} Evaluating a GP on gross IRR. Ignoring the timing advantage of deal-by-deal carry. Assuming fee offsets are standard.

\qffield{See also} \hyperref[def:alternative-fund-terms]{fund terms}; \hyperref[meth:private-fund-multiples-pme]{performance measures}.

\end{qfmethod}

\begin{qfmethod}[Private equity, private debt and special situations in a strategic allocation: comparing risk and return, roles, and a recommendation]\label{meth:private-capital-allocation}
\qffield{Statement} \emph{Return and risk across private capital.} Cambridge Associates statistics for US funds of the 1986--2008 vintages, as reported by Demaria and co-authors, risk-adjusted only through a PME \citep[pp.~116--117, 123, 130, 133]{demaria-private-markets-2021}:
\begin{center}\small
\begin{tabular}{@{}lrrrr@{}}
\toprule
Strategy & Pooled net IRR & Multiple & PME & Bottom/top quartile IRR \\
\midrule
Buyouts (604 funds) & 12.8\% & 1.71 & 1.41 & 6.5\% / 19.2\% \\
Distressed debt (110 funds) & 11.7\% & 1.56 & 1.36 & 5.5\% / 17.4\% \\
Growth equity (109 funds) & 19.3\% & 1.78 & 1.52 & 3.5\% / 19.1\% \\
Mezzanine (median fund) & 9.7\% & 1.48 & -- & 7.1\% / 11.9\% \\
\bottomrule
\end{tabular}
\end{center}
The growth-equity pooled IRR exceeds its top-quartile IRR as printed, which capital weighting by a few large funds could explain. Direct lending funds target 4--6 percent net in Western Europe and 7--9 percent gross in the US before fund leverage (p.~134).

\emph{Roles} (Interpretation, from the sources above and the private-equity, private-debt and diversification entries):
\begin{itemize}
\item \emph{Private equity} is a return-seeking substitute for public equity with leverage and manager risk; its premium over public equity is uncertain after risk adjustment \citep[pp.~427--428]{s7-ang}.
\item \emph{Private debt} provides contractual floating-rate income with covenants and lower dispersion, but credit risk that rises with leverage and the cycle; mezzanine sits between debt and equity.
\item \emph{Special situations} and distressed debt are countercyclical: their opportunity set widens in recessions, when buyouts and venture suffer \citep[p.~163]{demaria-private-markets-2021}.
\end{itemize}

\emph{Building the allocation.} Size the total private sleeve by liquidity, then split by role: private debt for income and lower dispersion, buyouts and growth for return, distressed as a cyclical complement. Diversify across managers and vintages, because manager dispersion exceeds strategy differences (Interpretation).

\qffield{Practice} \emph{A \$10 billion endowment} (original) with a 15 percent private sleeve and 4 percent annual spending: 7 percent buyouts and growth, 3 percent venture through established managers, 3 percent private debt (direct lending and some mezzanine) to fund near-term calls from income, 2 percent distressed with managers who can deploy in a downturn. On unsmoothed risk the sleeve behaves like leveraged equity, so the public equity weight is cut by more than the private equity added (\hyperref[meth:private-public-market-comparison]{private against public}).

\qffield{Conditions} Historical pooled returns are not expectations; they reflect survivorship in the databases and vintages with falling rates.

\qffield{Pitfalls} Comparing contractual yields on private debt with equity IRRs without the loss experience. Holding distressed debt for diversification but committing to it only after a crisis has begun.

\qffield{See also} \hyperref[meth:private-capital-diversification]{diversification}; \hyperref[meth:alternatives-allocation-approaches]{allocating to alternatives}.

\end{qfmethod}
